\documentclass[11pt,letterpaper]{article}
\usepackage[T1]{fontenc}
\usepackage[utf8]{inputenc}
\usepackage{lmodern,microtype}
\usepackage[margin=0.94in,headheight=15pt]{geometry}
\usepackage{amsmath,amssymb,amsthm,mathtools}
\usepackage{booktabs,longtable,tabularx,array}
\usepackage{xcolor,listings,enumitem,needspace,etoolbox,fancyhdr}
\usepackage{tocloft}
\setlength{\cftbeforesecskip}{0pt}
\usepackage{xurl,hyperref,bookmark}
\definecolor{heather}{HTML}{653E70}
\definecolor{ink}{HTML}{26313C}
\definecolor{muted}{HTML}{626977}
\definecolor{wash}{HTML}{F5F1F7}
\hypersetup{colorlinks=true,linkcolor=heather,citecolor=heather,urlcolor=heather,
 pdftitle={Heather: Proof-Carrying Mathematical Interfaces},
 pdfauthor={ChatGPT, prepared for Vladimir Reshetnikov},
 pdfsubject={Round-four proposal for property- and behavior-aware mathematical elaboration over Lean}}
\urlstyle{same}
\pagestyle{fancy}\fancyhf{}
\fancyhead[L]{\small\sffamily\color{heather}HEATHER}
\fancyhead[R]{\small\sffamily Proof-carrying mathematical interfaces}
\fancyfoot[C]{\small\thepage}
\setlength{\parindent}{0pt}\setlength{\parskip}{0.45em}
\setlength{\emergencystretch}{2.5em}
\setcounter{tocdepth}{1}
\setlist{itemsep=0.22em,topsep=0.35em,leftmargin=*}
\newcolumntype{Y}{>{\raggedright\arraybackslash}X}
\newcolumntype{P}[1]{>{\raggedright\arraybackslash}p{#1}}
\newcommand{\code}[1]{\texttt{\detokenize{#1}}}
\newcommand{\N}{\mathbb N}\newcommand{\Z}{\mathbb Z}\newcommand{\Q}{\mathbb Q}\newcommand{\R}{\mathbb R}
\newcommand{\Beh}{\mathsf{Beh}}\newcommand{\Req}{\mathsf{Req}}
\newcommand{\Spec}{\mathsf{Spec}}\newcommand{\ev}{\mathsf{eval}}
\newcommand{\obs}{\mathsf{obs}}\newcommand{\id}{\mathsf{id}}
\newcommand{\erase}[1]{\lvert#1\rvert}
\newcommand{\typing}{\vdash}
\newtheorem{theorem}{Theorem}[section]
\newtheorem{proposition}[theorem]{Proposition}
\newtheorem{lemma}[theorem]{Lemma}
\newtheorem{corollary}[theorem]{Corollary}
\theoremstyle{definition}\newtheorem{definition}[theorem]{Definition}
\theoremstyle{remark}\newtheorem{remark}[theorem]{Remark}
\lstset{basicstyle=\small\ttfamily,columns=fullflexible,keepspaces=true,
 breaklines=true,showstringspaces=false,backgroundcolor=\color{wash},
 frame=single,rulecolor=\color{wash},framesep=6pt,
 xleftmargin=6pt,xrightmargin=6pt,aboveskip=0.7em,belowskip=0.7em}
\pretocmd{\section}{\Needspace{11\baselineskip}}{}{}
\pretocmd{\subsection}{\Needspace{6\baselineskip}}{}{}
\begin{document}
\hypersetup{pageanchor=false}
\begin{titlepage}
\vspace*{12mm}
{\sffamily\large\color{heather}MATHEMATICAL PROOF-LANGUAGE DESIGN\par}
\vspace{18mm}
{\sffamily\bfseries\color{heather}\fontsize{46}{50}\selectfont Heather\par}
\vspace{8mm}
{\sffamily\bfseries\fontsize{25}{30}\selectfont Proof-Carrying\\Mathematical Interfaces\par}
\vspace{8mm}
{\sffamily\Large A contract-directed language over Lean\par}
\vspace{13mm}
\hrule\vspace{5mm}
\textbf{Fourth-round thesis.} A mathematical object should bring its proved
properties to every legitimate use, without changing its identity. An operation
should expose the precise evidence it needs. Automation should connect the two,
while preserving the statement, the intended argument, and the scope of every
assumption.
\vspace{7mm}
Heather combines evidence-implicit parameters, snapshot-local elaboration,
relational behavioral contracts, and domain-sensitive finite certificates.
The contribution is a concrete integration design and an executed narrow
reference model, not another proposal to put unrestricted theorem proving
inside type conversion.
\vfill
Prepared for Vladimir Reshetnikov\par
6 September 2026\par
\vspace{5mm}
{\small\color{muted}Evidence status: both round-three synthesis reports and selected
primary sources were reviewed. Eighteen Python test methods pass. Generated Lean
proof candidates are included but were not compiled in this environment.
No full repository build, verified frontend, Leant integration, or productivity
measurement is claimed.}
\end{titlepage}
\hypersetup{pageanchor=true}\pagenumbering{arabic}

\begin{abstract}
Heather is a proposed mathematical proof language whose central unit is a
\emph{proof-carrying interface to an unchanged object}. Rather than asking
mathematicians to repeatedly supply already-known side conditions, Heather
makes selected proof parameters implicit at use sites and reconstructs their
arguments from a bounded, theorem-backed evidence index. Meaning-bearing data
remain explicit or uniquely determined before proof search: the chosen ring,
topology, measure, region, witness specification, and quantifier structure are
not variables that search may adjust to obtain an easier theorem.

This fourth-round proposal accepts the main conclusion of the two round-three
syntheses: the first implementation should extend Lean's author-facing typing
and elaboration discipline, not its trusted conversion relation. It advances
that conclusion in three directions. First, it specifies an evidence-implicit
interface with concrete scoping, rewriting, rule-generation, and publication
contracts. Second, it develops \emph{realizable affine frames}: domain-sensitive
finite certificates that decide affine observation equalities on the actual
admissible input space, including empty element types and correlated inputs.
A rank argument identifies the number of observations needed for unrestricted
affine equality. Third, it supplies an executed Python frontend for a restricted
list language, a certificate rechecker, concrete negative witnesses, a small
use-site requirement resolver, and ordinary Lean exports.

The article distinguishes mathematical proofs, executed finite tests, historical
kernel checks reported by the source documents, and new kernel checks, of which
none were obtained here. It works through the ProveIt autonomous-derivative
argument, the Fermat repository's finite tensor representability construction,
and the inherited exact-quotient benchmark. It also gives an implementation
contract for Leant and certified computer algebra, evaluates genuine kernel
extension alternatives, and states stopping rules for the separate-language
experiment.
\end{abstract}

\begingroup\small\setlength{\parskip}{0pt}\tableofcontents\endgroup
\clearpage

\section{The decision: hide routine evidence, not mathematical choices}

A mathematician who has introduced a positive real number does not normally
write a separate proof of nonzeroness before every division. A mathematician
who differentiates an identity on an open set nevertheless relies on a real
locality argument: the identity holds near the point, not merely at that point.
These two omissions are similar in prose but different in implementation. The
first needs a small implication. The second needs the right region, topology,
point, and quantified hypothesis before a derivative congruence theorem applies.

Heather aims to make both steps concise \emph{without making either step
unaccountable}. Its design principle is:
\begin{quote}
Keep the mathematical nouns, choices, and argument landmarks in the source.
Infer routine evidence at their interfaces. Retain an ordinary proof of every
assertion that the document publishes.
\end{quote}
This does not promise that every intuitive mathematical step becomes automatic.
It promises that repeated interface work can be delegated under a precise,
inspectable contract, and that a blocked operation reports the mathematical
premise it still lacks rather than merely exposing a tactic failure.

\subsection{What Heather takes from round three}
The reports \emph{From Properties to Proof Obligations} and \emph{Nine
Refinements} are closely aligned on the foundational choice. Both distinguish
properties of a fixed object from choices of structure, use ordinary quantified
propositions for behavior, require proof-producing inference, and recommend
building an actual use-site interface before adding a large narrative syntax.
Both also warn that the proposed interface must beat ordinary Lean supplied
with the same new libraries and services \cite{round3a,round3b}.

Heather adopts that architecture. It does not rename the four standard
refinement rules and call their agreement a new discovery. Nor does it count
historical compilation of small contract combinators as implementation of a
language. The source reports themselves distinguish these levels of evidence,
and their reconciled versions correct important overstatements about finite
observations, \code{List Empty}, theorem counts, and the meaning of compilation
\cite{round3a,round3b}.

\subsection{The proposed increment}
The intended fourth-round increment is summarized below. ``Executed'' refers
to this archive's Python reference model, not to Lean acceptance.

\begin{center}\small
\begin{tabularx}{\textwidth}{@{}P{0.23\textwidth}Y P{0.20\textwidth}@{}}
\toprule Question & Heather's answer & Present evidence\\\midrule
What should a use site do? & Insert only designated proof arguments under a fixed
meaning; expose missing premises and their consumers. & Specified; tiny Nat
requirement resolver executed.\\
What survives rewriting and edits? & Exact proof anchors within one snapshot;
checked transport within a context; rebuild after declaration edits by default.
& Specified; cross-anchor and stale-epoch refusals executed.\\
When can finite observations decide behavior? & When a semantic model and a
realizable, covering frame determine the demanded relation on the admitted domain.
& Written theorems; four list-domain instances executed.\\
What is the kernel extension? & An author-facing typing discipline first;
a primitive refinement former only after a measured encoding bottleneck.
& Design and translation argument.\\
What would justify the language? & Lower authoring or repair cost than equally
equipped Lean, without increased semantic-reading errors. & Evaluation design;
no human study.\\
\bottomrule
\end{tabularx}
\end{center}

The affine-frame result is a useful mathematical interface theorem, not a claim
of historical novelty in linear algebra. Its contribution to this campaign is
to turn the empty-type and correlated-observation warnings into a positive,
domain-adaptive acceptance method with a precise completeness statement.

\section{Source review and a fair Lean baseline}\label{sec:sources}

\subsection{Snapshots and scope}
The repositories were inspected through their GitHub interfaces. The relevant
snapshot identities are retained in Appendix~\ref{app:sources}. ProveIt and the
Fermat repository resolve to the same commits discussed by the round-three
reports. The current Leant \code{main} returned in this review resolves to
\code{3a40904b}; the syntheses additionally discuss the distinct snapshot
\code{823259f7}. Heather does not silently identify those two versions. Its direct
Leant source observations below concern \code{3a40904b}; claims about the other
snapshot remain attributed to the reports \cite{leant-verification,leant-adapter,
round3a,round3b}.

The review is selective rather than a repository audit. It closely studies the
two requested synthesis texts and representative implementation files. It does
not claim to have rebuilt ProveIt or the Fermat development, proved the whole
Fermat theorem here, reviewed every round-three proposal independently, or
measured the distribution of verbosity throughout either repository.

\subsection{ProveIt: a genuinely local induction argument}
The file \path{AutonomousIteratedDeriv.lean} proves a reusable theorem about
an autonomous differential equation. On an open set $U$, suppose
\[
 W'(x)=\varphi(W(x)),\qquad G_0(W(x))=W(x),
\]
and the derivative of $G_n$ exists at each relevant $W(x)$ with a supplied
value $G'_n(W(x))$. If
\[
 G_{n+1}(W(x))=G'_n(W(x))\,\varphi(W(x)),
\]
then $D^nW(x)=G_n(W(x))$ for every $n$ and $x\in U$
\cite{autonomous}.

The Lean proof performs induction before introducing the point in each branch.
In the successor step, it turns the induction hypothesis on $U$ into eventual
equality at the chosen point, uses derivative congruence, applies the chain
rule, and rewrites by the recurrence. This is not purposeless verbosity: each
operation has a mathematical role. The avoidable cost is repeatedly connecting
those roles and their parameters by hand. Equally important, the theorem asks
for differentiability of $G_n$ only on $W(U)$, not on all of $\R$. A proposed
shortening that assumes global differentiability has changed the theorem
\cite{autonomous}.

The same file already contains a concise corollary obtained by applying this
reusable theorem to globally differentiable $G_n$. Heather should preserve that
library-level compression. Once a theorem packages the right mathematics,
a one-line application is a good outcome, not an embarrassment for the new
language.

\subsection{Fermat: structure assembly and a relational guarantee}
The file \path{P2M/Sol/S_Algebra_exists_algHom_equiv_pi.lean} constructs a finite
free representing algebra for a finite family of finite free commutative
$A$-algebras. Its proof separates a two-factor tensor universal property from
induction over a finite index type. It assembles structure instances,
equivalences of homomorphism spaces, and a naturality equation
\cite{tensor}.

This example has three distinct costs. There is mathematical construction:
use the base algebra for the empty family and a tensor product for a new factor.
There is interface plumbing: instantiate module facts and compose equivalences.
Finally there is a relational proof: the chosen equivalence commutes with
postcomposition. Heather should reduce the second cost and expose the first
and third. Naturality is not an adjective attached independently to each
component map.

The original theorem explicitly assumes that \emph{every factor} is finite and
free as an $A$-module. The round-three syntheses correctly identify an exposition
that lost those hypotheses. Heather's type layer must fail to infer module
finiteness when a factor's finiteness is removed; it must not infer it from
finiteness of the index type alone \cite{tensor,round3a,round3b}.

\subsection{Leant: association is already treated seriously}
Leant's \path{Synth/Verification.hs} makes \code{Verified} an opaque,
nominally parameterized callback-acceptance receipt. Its documented boundary
explicitly distinguishes this receipt from a behavioral theorem, solver
certificate, or kernel proof. Candidate variants are tried in order, and only
accepted groups consume the success quota \cite{leant-verification}.

The \path{PostVerification.hs} module scopes candidate occurrences by a
rank-two epoch and validates a proposed permutation against the original
receipts. This is useful association machinery: it protects which occurrence
was accepted and subsequently selected. It does not, by itself, prove the
candidate's mathematics \cite{leant-post}.

The Length adapter is similarly explicit. It ties an accepted candidate graph
to a sealed query and keeps the product-domain path distinct. Its comments
warn that replay of an abstract finite-spine counterexample does not authorize
a source-level behavioral refutation or candidate pruning. Heather therefore
proposes an additional semantic bridge, not a repair for a bug that the source
already disclaims \cite{leant-adapter}.

\subsection{Existing facilities are part of the comparison}
Lean already has extensible elaboration and proof-producing tactics. Its
function-application machinery includes automatic parameters whose associated
tactic scripts construct omitted arguments. Thus the novelty of Heather's
\code{requires} interface cannot be merely ``a tactic fills a proof parameter''
\cite{lean-reference,lean-application}.

Mathlib's \code{fun_prop} already composes registered function-property theorems.
Lean's \code{mvcgen} supports verification-condition generation from program
specifications. Library search, simplification, arithmetic automation, and
appropriate uses of \code{grind} are also legitimate controls
\cite{funprop,mvcgen,lean-tactics}. The report-reviewed 4.32.0 supply probe
has version-specific qualifications, including experimental \code{mvcgen}
support; the current online manual is not the same snapshot as ProveIt's
pinned toolchain \cite{round3a,round3b,lean-reference}.

Heather's hypothesis is an integration hypothesis: a common interface for
requirements, fixed meaning, explanation, scope, and replay may be more useful
than independent tactics at independent holes. It must earn that claim by
comparison, not by pretending the host lacks automation.

\section{What the mathematician would write}\label{sec:surface}

\subsection{A small vocabulary with different logical effects}
Heather reserves different verbs for different commitments. \code{given}
introduces data; \code{assume} introduces a hypothesis into the stated theorem;
\code{establish} requests a proof; \code{choose} authorizes construction of a new
witness satisfying a fixed specification; \code{observe} requests a specified
view of an existing object. \code{use} applies a registered mathematical method.
The distinctions matter more than these particular English spellings.

For example, the following is \emph{proposed full Heather syntax}, not the
implemented four-line grammar of the companion:
\par\noindent\begin{minipage}{\linewidth}
\begin{lstlisting}
given x : Real
assume hx : 0 < x
establish reciprocal_cancels : x * (1 / x) = 1
  via field cancellation
\end{lstlisting}
\end{minipage}\par
The author supplies positivity, the target, and the method. The field
cancellation interface requests a proof of $x\ne0$; the evidence engine derives
it from \code{hx}. The generated proof parameter is hidden in the compact view,
not absent from the accepted proof. If the assumption is weakened to $0\le x$,
the nonzero obligation remains open. Heather does not silently add $x\ne0$ to
the theorem's hypotheses.

An assumption introduced during a proof changes its statement unless it is
inside an explicit conditional subproof. Consequently an editor command that
repairs a failed proof by adding an assumption must present a statement change,
not report that it has merely repaired elaboration.

\subsection{Keep induction general enough}
For the autonomous equation, a compact mathematical argument should resemble:
\par\noindent\begin{minipage}{\linewidth}
\begin{lstlisting}
establish derivatives :
  for every n : Nat, for every x in U,
    D[n](W)(x) = G[n](W(x))
proof
  induct n, retaining the quantifier over x in U
  zero:
    conclude by the initial equation
  successor n:
    on U, D[n](W) = G[n] compose W by induction
    differentiate this identity locally on U
    conclude by the chain rule and the recurrence
\end{lstlisting}
\end{minipage}\par
The successor line represents a method application with several proof
parameters, not an unrestricted natural-language instruction. Its interface
requires equality on a neighborhood of the point. The open-set evidence and
membership supply that neighborhood; the derivative hypotheses on the actual
range supply the chain-rule inputs. Section~\ref{sec:analysis} gives the full
mathematical proof and the obligations represented by these lines.

The compact source must preserve the quantifier over $x$. A point-specialized
induction hypothesis does not suffice for differentiating an identity. Heather
therefore makes the retained induction motive part of the elaborated method
interface, rather than relying on a language model to guess the missing scope.

\subsection{Reason as hint, or reason as obligation}
Two annotations have intentionally different meanings:
\par\noindent\begin{minipage}{\linewidth}
\begin{lstlisting}
establish P  hint by induction
establish P  via induction on n
\end{lstlisting}
\end{minipage}\par
The first influences search and permits any accepted proof of $P$. The second
requests a proof represented by an induction method node with checked base and
step obligations. It is insufficient for the emitted Lean term to contain an
unused reference to an induction theorem. Heather validates the designated
method's derivation interface, including its conclusion and children.

This mechanism checks an explicit operational interpretation of ``by
induction'', not a human author's private intention. A method registry must say
what each mandatory method requires. Some methods may be deliberately broad;
others may restrict permissible support or require a particular transformation.
The policy is visible in the document.

\subsection{Two views, one checked document}
The compact view displays theorem statements, mathematical choices, significant
intermediate assertions, open conditions, and required methods. The expanded
view displays resolved parameters, selected structures, evidence dependencies,
search routes, and exported Lean terms. Both views are projections of the same
typed document model.

A statement such as ``the solution set is complete'' cannot be shortened to
``a solution exists'' in the expanded target, even if the latter is easy to
prove. Likewise, a published intermediate assertion must have its own accepted
proof even if it is not used by the final theorem. Checking only the final proof
would allow a document to contain unchecked mathematics in its explanatory
middle. These publication rules are inherited from the earlier campaign and
retained by both round-three syntheses \cite{round3a,round3b}.

\section{An evidence-aware type discipline}\label{sec:types}

\subsection{The judgment and its ordinary Lean meaning}
A Heather typing state contains an immutable declaration environment $\Sigma$,
a dependent local context $\Gamma$, a set of selected structure arguments, and
an evidence index $I$. Structure arguments are ordinary data within the
elaborated context; listing them separately is an implementation and explanation
convenience, not a second logic.

For a fixed Lean term $a:A$, a \emph{facet} is a proposition $P(a)$ together
with a retained proof $p:P(a)$. A row $\Phi$ is a finite collection of such
propositions. It may mention several objects: $\operatorname{Commutes}(f,g)$,
$\operatorname{EqOn}(f,g,U)$, or a naturality law for a whole family. The useful
source judgment has the form
\[
 \Sigma;\Gamma;I\typing t \Rightarrow A\mid\Phi
 \quad\leadsto\quad a:A,\quad \pi:\bigwedge_{P\in\Phi}P(a).
\]
Here $A\mid\Phi$ is an author-facing classification. Locally, Heather does
\emph{not} replace $a$ by a new subtype value every time a fact is learned.
Instead, it records ordinary Lean proofs in $I$. An explicit boundary can
package $a$ and selected proofs into a subtype or structure when an API actually
expects such a value.

A row is not a collection of informal tags. In particular, \code{positive}
has a fully elaborated proposition containing its ordered structure, zero,
and carrier. \code{integrable} contains its measure, measurable structures,
codomain, and function. No matching on the printed word alone is meaningful.
The metadata may index a tuple of subjects for efficient retrieval, but the
proof's type is authoritative.

\subsection{Four elementary rules and two non-elementary obligations}
The familiar rules are retained, with their proof terms made explicit:
\[
\frac{a:A\quad p:P(a)}{a:A\mid P}
\qquad
\frac{p:P(a)\quad q:Q(a)}{a:A\mid P\wedge Q}
\]
\[
\frac{p:P(a)\quad h:\forall x,\,P(x)\to Q(x)}{a:A\mid Q}
\qquad
\frac{p:P(a)\quad e:a=b}{b:A\mid P}.
\]
Their translations use the supplied proof, conjunction introduction, $h(a,p)$,
and equality elimination. There is no semantic subtyping search hidden inside
Lean's conversion checker. Subsumption is a generated proof argument.

Two less elementary obligations deserve equal prominence. First, changing an
\emph{object} requires a relation to the original object or an explicitly
specified construction; row conjunction cannot identify two separately chosen
witnesses. Second, changing a \emph{context} requires a well-typed substitution
for its dependent telescope. A matching printed name, a rewritten string, or a
matching cache digest is not such a substitution.

\subsection{Property-implicit parameters}
Heather introduces an argument role distinct from both ordinary data inference
and typeclass search:
\par\noindent\begin{minipage}{\linewidth}
\begin{lstlisting}
method cancel_left (x y z : R)
  requires (hx : Regular x)
  requires (hxy : x * y = x * z)
  ensures y = z
\end{lstlisting}
\end{minipage}\par
The declaration denotes an ordinary theorem with explicit proof parameters.
At a Heather call site, designated \code{requires} arguments become goals for
the evidence engine. Explicitly supplied proofs always remain available as an
escape hatch. The implementation should initially realize this role through
Lean metaprogramming and automatic-parameter machinery, plus registry metadata,
rather than by modifying the kernel \cite{lean-application}.

Only a proposition-valued parameter may have this role. A ring structure,
measurable space, enumeration, inverse function, or normalization witness is
not a proof argument merely because constructing it is inconvenient. Such data
are inferred by the ordinary elaborator when uniquely determined, supplied by
the author, or requested through an explicit construction specification.

A method may also have dependent output data. Its semantic shape is then
\[
 (x:A)\to (h:P(x))\to\{y:B(x)\mid Q(x,y)\}.
\]
The result creates a new $y$, not a new view of $x$. Later requirements may
mention that exact $y$ and proofs returned with it. An unordered set of guard
strings cannot substitute for this dependent binding order.

\subsection{Refinement binders are explicit about quantifiers}
A source binder
\par\noindent\begin{minipage}{\linewidth}
\begin{lstlisting}
given x : A satisfying P(x)
\end{lstlisting}
\end{minipage}\par
introduces $x:A$ and a proof of $P(x)$ as assumptions of the surrounding
statement. Accordingly,
\[
 (\forall x:A\text{ satisfying }P(x),\ Q(x))
 \quad\rightsquigarrow\quad \forall x:A,\ P(x)\to Q(x),
\]
whereas existential syntax translates to
$\exists x:A,\ P(x)\wedge Q(x)$. The word \code{establish} never has this
assumption-introducing effect.

The surface deliberately avoids an ambiguous bare arrow such as
``$A\mid P\to B\mid Q$''. Two legitimate readings are different:
\[
 f:A\to B\quad\text{with }\forall x,\ P(x)\to Q(x,f(x)),
\]
and
\[
 f:(x:A)\to P(x)\to B.
\]
The former describes a total function with a guarded behavioral guarantee.
The latter only provides an evaluation interface when a guard proof is
available. Heather requires these two forms to be spelled differently.
It never silently extends a refined-domain function to all of $A$.

\subsection{A meaning boundary, not a hash-based proof rule}\label{sec:meaning}
Before invoking external search, Heather records the elaborated target and its
meaning-bearing arguments: carrier types, universe instantiations, chosen
structures, fixed objects, quantifiers, relations, regions, measures, precision,
and permitted result strength. Proof holes may remain; an explicitly authorized
witness hole may also remain under its fixed specification. Other unsolved data
holes are resolved or reported as ambiguities before the search is accepted.

This policy prevents a prover from satisfying a request by choosing the zero
ring, a different measure, a weaker postcondition, or a convenient branch of
an inverse function that the author did not request. It does not prevent
ordinary mathematical construction: a request to choose \emph{any} root meeting
a stated condition really does authorize a choice among such roots.

The implementation may compute a digest of this record for routing or caching.
Acceptance still checks the actual term against the actual target in the
actual environment. A digest establishes neither equality of mathematical
objects nor a theorem about their behavior. The reference companion follows
this separation: the request digest is checked, but certificate data are also
reconstructed from the original parsed request.

\subsection{Restricted preservation theorem}
\begin{theorem}[Conservative interpretation of the core rules]\label{thm:preserve}
Fix a Lean environment and a well-typed interpretation of the source's data,
structures, and statement. Suppose a Heather derivation uses only ordinary
term formation, the four facet rules above, designated proof-parameter
application, and explicit dependent constructions. Suppose every external
proof leaf is accepted at its fixed Lean proposition, and every registered
rule is applied with well-typed arguments. Then the interpreted derivation
produces a Lean term of the interpreted source statement, with no additional
axioms beyond those used by its leaves and registered theorems.
\end{theorem}
\begin{proof}
Induct over the derivation. Ordinary term rules are interpreted by the
corresponding Lean constructors. Facet introduction retains a proof; conjunction
uses its constructor; consequence applies the registered implication; transport
uses equality elimination. A proof-implicit call inserts the recursively
constructed proof into an ordinary dependent function application. A
construction introduces its actual output and associated proofs into the
translated dependent context. Each rule therefore preserves the interpreted
target under the translated context. Since no rule introduces an axiom, the
axiom dependencies are inherited from the terms used at its leaves and
applications. At the root, the result has the fixed interpreted source type.
\end{proof}

The hypothesis about source interpretation is essential. A buggy parser could
interpret the wrong quantifier order and still generate a perfectly valid Lean
proof of the wrong statement. This theorem does not verify the parser, registry
loader, editor, renderer, or Python implementation. It is the mathematical
contract those components must eventually implement and test.

\section{Behavior is part of an interface, not a confidence label}\label{sec:behavior}

\subsection{Unary and relational contracts}
For a fixed total function $f:A\to B$, define
\[
 \Beh(P,Q,f)\;:\!\!\iff\;\forall x:A,\ P(x)\to Q(x,f(x)).
\]
The postcondition is relational: it mentions the input and the actual output.
An output-only predicate cannot express preservation of a given list's
multiset, a chosen element's orbit, or a computation's relation to its original
expression.

Other laws quantify over multiple calls. Monotonicity, Lipschitz continuity,
commutation, naturality, and prefix locality are ordinary predicates about the
whole function or configuration. Heather permits them as facets without
pretending that each can be reduced to a unary postcondition. For example,
\[
 \operatorname{Lipschitz}(L,f)\iff
 \forall x,y,\ d(f(x),f(y))\le L\,d(x,y)
\]
retains the two metric structures and the value of $L$. A proof of this
property is an upper-bound certificate, not automatically a proof that $L$ is
the least possible constant.

\subsection{Consequence with the correct directions}
Suppose $\Beh(P,Q,f)$, and suppose
\[
 \forall x,\ P'(x)\to P(x),\qquad
 \forall x,y,\ P'(x)\to Q(x,y)\to Q'(x,y).
\]
Then $\Beh(P',Q',f)$: apply the first implication to obtain the old input guard,
use the original contract, then weaken its postcondition using the second
implication. This is the familiar contravariance of preconditions and covariance
of guarantees. Calling a contract simply ``stronger'' obscures which side is
being changed.

Heather stores implication theorems rather than an untyped ranking of contracts.
Equality, equality on a region, almost-everywhere equality, interval enclosure,
and a finite jet are different relations, not confidence scores. Some
promotions are possible under additional hypotheses; their exact theorems
belong to consumer interfaces.

\subsection{Composition retains the intermediate}
Let $f:A\to B$ and $g:B\to C$. Assume
\[
 \Beh(P,Q,f),\qquad \Beh(R,S,g),
\]
with bridges
\begin{align*}
 &\forall x,y,\ P(x)\to Q(x,y)\to R(y),\\
 &\forall x,y,z,\ P(x)\to Q(x,y)\to S(y,z)\to T(x,z).
\end{align*}
Then $\Beh(P,T,g\circ f)$. Indeed, at an input $x$ satisfying $P$, retain
$y=f(x)$, obtain $Q(x,y)$, establish $R(y)$, obtain $S(y,g(y))$, and use the
second bridge. Matching the carrier types $A\to B\to C$ does not prove either
bridge. If $g$ also depends on $x$ or on earlier constructed data, use the
corresponding dependent version and preserve that binding order.

A requirement transformer for a fixed $f$ is a sufficient condition generator
with a theorem
\[
 \Req_f(Q)(x)\Longrightarrow Q(f(x)).
\]
Composition gives the sufficient transformer
$\Req_f(\Req_g(Q))$ for $g\circ f$. Heather does not claim it is a weakest
precondition over arbitrary mathematical predicates. Two registered sufficient
routes can be incomparable; choosing one need not solve a universal optimization
problem over all proofs.

\subsection{Sorting and higher-order obligations}
A request for $\mathsf{List}(A)\to\mathsf{List}(A)$ admits a constant-empty
function. Requiring length preservation still admits operations that discard
all input values and repeat a fixed value. Sortedness alone also admits
constant empty output. A sorting specification therefore needs a relation to
the actual input, such as permutation, together with the required ordering law.
Stability, when requested, is a further property rather than an implication of
those two clauses.

For a higher-order client, the same discipline applies. A combinator accepting
an order-preserving function must receive a proof about the \emph{function
passed to it}, under the same orders. A theorem about another extensionally
equal function can be used after a checked transport; a convenient function
name or a nominal provider tag cannot stand in for that transport.

A polymorphic Lean type does not automatically supply every familiar
parametricity theorem. Classical constructions can inspect properties of a
type; the Lean reference gives an explicit example. Heather therefore requires
an actual uniformity theorem or a restricted syntactic fragment with a proved
semantic result before it transfers behavior across element types
\cite{lean-axioms}.

\subsection{Effects are a separate dimension}
A mathematical property of a pure total function, partial correctness of a
stateful program, termination, and a search timeout are distinct concepts.
For effectful code, Heather should reuse a suitable Hoare-logic interface and
Lean's verification-condition infrastructure rather than relabeling every
behavioral predicate as an effect \cite{mvcgen}.

An external solver may time out while a theorem remains true. Conversely, an
algorithm may return quickly but fail its postcondition. Heather records
operational budgets in the request protocol and logical premises in the proof
context. A timeout is never inserted as a mathematical assumption.

\section{Bounded inference and the obligation frontier}\label{sec:inference}

\subsection{A registry is a typed mathematical API}
A registered rule records an existing theorem, roles for its parameters,
which premises are eligible for property inference, the conclusion pattern,
and a selected inference profile. Registration is validated against the
actual theorem type. Metadata cannot authorize a conclusion the theorem does
not have.

There should be no global instance for every derived fact. Instances select
structure or serve existing APIs; the evidence index supplies proof arguments.
When an imported theorem genuinely expects a proposition packaged as an
instance, a local, tightly scoped adapter can provide it. This avoids confusing
ordinary forward reasoning with open-ended typeclass search.

The index need not be large. Most useful rules should either discharge a
frequent missing interface condition or connect facts from two genuinely
different packages. A rule that duplicates cheap lookup but increases search
branching should be removed unless it improves explanation or robustness.

\subsection{How the finite ground fragment is generated}
The finite-closure argument of the round-three syntheses starts after a ground
rule set has been selected. Heather specifies a conservative first generator
rather than treating that step as free \cite{round3a,round3b}.

Start from a finite typed pool $V$ containing elaborated subterms of the target,
selected local declarations and their types, and explicitly named
intermediates. A profile may generate additional terms using a finite operator
list, but only up to a declared depth and a total term limit. It may not keep
inventing larger terms until something works. Data parameters of a rule are
instantiated from compatible entries in this fixed pool; proof parameters
become premises. Dependent parameter types are checked in order. Universe and
structure arguments must match the fixed elaborated terms.

If the resulting pool has $M$ entries and rule $r$ has $a_r$ eligible data
parameters, the raw number of candidate substitutions is bounded by
$M^{a_r}$. Thus $\sum_r M^{a_r}$ is a crude finite upper bound before type
filtering. It is not a promise that this generation is cheap. Indexes on types,
heads, and roles should reduce the practical enumeration. Each profile also
bounds the number of attempts and stops with a distinct generation-exhausted
status when its limits are reached.

After constructing the candidate rules, take the backward dependency slice of
the demanded goals, including the premises of every retained rule. Forward
closure on this fixed slice is predictable. A theorem requiring a witness or
intermediate term outside the pool may be missed. That is a limitation of the
profile, not evidence against the theorem. The author can introduce the
intermediate, choose another profile, supply a proof, or invoke a separately
budgeted provider.

\subsection{What the closure theorem really establishes}
Let $\mathcal Q$ be the finite set of propositions in the chosen fragment and
let $\mathcal R$ contain rules
\[
 P_1\wedge\cdots\wedge P_m\longrightarrow Q
 \quad(P_i,Q\in\mathcal Q),
\]
each instantiated from a theorem. Starting from accepted evidence $S_0$, add a
conclusion only after all its premises have proofs.

\begin{proposition}[Finite evidence closure]
A fair worklist procedure adds at most $|\mathcal Q\setminus S_0|$ new facts,
retains an acyclic derivation for each added fact, and computes the least
$\mathcal R$-closed superset of $S_0$. Different fair schedules compute the
same set of facts, though they may retain different proofs.
\end{proposition}
\begin{proof}
Each addition is new and belongs to the finite set $\mathcal Q$, proving the
bound. Its premises were established earlier, so selecting their existing
proofs yields an acyclic derivation and a proof of the conclusion. Every closed
superset of $S_0$ contains each addition by induction over the worklist. Fairness
ensures that every enabled rule is eventually considered, so the final set is
itself closed. Leastness and schedule independence follow.
\end{proof}

With premise counters, worklist bookkeeping can be linear in the number of
facts and premise occurrences. Rule generation, dependent unification,
normalization, term construction, proof size, and kernel checking are excluded
from that statement. Heather makes no principal-refinement theorem or global
completeness claim for Lean.

\subsection{A frontier is a useful partial result}
When a demand is not discharged, Heather records the proposition, local
context, consuming operation, and a small set of available sufficient routes.
For example:
\par\noindent\begin{minipage}{\linewidth}
\begin{lstlisting}
Operation: cast the exact quotient into K
Resolved: 4 divides numerator in Int
Open:     the image of 4 is a unit in K
Route:    characteristic zero and field structure suffice
Note:     nonzero in Int is not this destination condition
\end{lstlisting}
\end{minipage}\par
This is not a proof that no other route exists. It is a valid dependency slice
for the attempted methods. Finding a globally smallest explanation can itself
be expensive; the first implementation should prefer correct and stable
explanations over unproved minimality.

An unseeded cycle, such as $P\to Q$ and $Q\to P$, contributes no fact. A
conditional result whose guard is $G$ may not be used as unconditional evidence
to prove $G$. A real induction or fixed-point argument must enter through its
own theorem-backed method rather than through a cyclic cache entry.

\section{Evidence through rewriting, branches, and edits}\label{sec:lifetime}

\subsection{The authoritative entry is a typed proof}
An implementation entry has the conceptual form
\[
 (\kappa,\Gamma_0, P, p, \text{origin}, \text{dependencies}),
 \qquad\Gamma_0\typing p:P,
\]
where $\kappa$ identifies an elaboration snapshot. Indices on expression heads,
subjects, or normal forms accelerate retrieval, but do not change this invariant.
A proof can be offered to a consumer only in a context where its dependencies
are valid and its type matches the requested proposition, possibly through an
explicit checked adapter.

This design deliberately reuses Lean's contexts and terms. Heather adds an
index and explanations; it does not introduce an independent database of
mathematical truth whose agreement with Lean is assumed.

\subsection{Three cases after simplification}
Suppose $p:P(t)$ is available and simplification replaces a displayed expression
$t$ by $u$.

If $t$ and $u$ are definitionally equal in the fixed context, ordinary Lean
checking can reuse the proof. If a checked equality $e:t=u$ is available,
equality elimination transports $p$ to $P(u)$. If only a weaker relation is
known, such as equality almost everywhere, a jet relation, or an order bound,
then a consumer-specific theorem must justify the requested consequence.

An equivalence of propositions is also useful: from
$h:P(t)\leftrightarrow Q(u)$ and $p:P(t)$, the forward implication yields
$Q(u)$. This does not assert $t=u$. The transport trace records which operation
occurred, so the expanded view does not advertise equality of objects when only
equivalence of their relevant properties was used.

For dependent indices, the predicate family and its type may need transport
together. A term indexed by $n$ cannot be treated as one indexed by $m$ merely
because a pretty-printer suppresses the index. The first native prototype
should support a few explicit equality and proposition-equivalence adapters,
not a general coercion planner.

\subsection{Context morphisms and branch export}
Let $\theta$ be a well-typed substitution interpreting every declaration of
$\Gamma_0$ in a destination context $\Delta$. It must respect the dependent
order: the interpretation of a later declaration is checked after substitution
in its type. The ordinary substitution principle gives
\[
 \Gamma_0\typing p:P
 \quad\Longrightarrow\quad
 \Delta\typing p[\theta]:P[\theta].
\]
This is the legitimate basis for cross-context reuse. In a native implementation,
the destination Lean checker validates the transported term; the context map
is not justified by a correspondence between variable names.

A branch adding $h:H$ may use $h$ locally. On leaving the branch, its result
$P$ exports as $H\to P$, or under an appropriate case-analysis theorem, not as
an unconditional fact for a sibling branch. The same rule applies to an
existential witness introduced in a local proof: its scope and dependencies
cannot be replaced by a global name with a similar printed form.

\subsection{Rebuild first, migrate later}
Heather's default after a declaration edit is to discard its local evidence
index and reconstruct it from the newly elaborated context. Long-lived caches
may retain candidate theorem names, source edits, or proof artifacts for replay,
but not a claim that an old local identifier is valid in the rebuilt declaration.
A successful replay produces a new acceptance record.

This is intentionally conservative. A persistent proof cache with verified
context embeddings may eventually save work, but it introduces a substantial
engineering surface. The first prototype should measure rebuild cost before
making migration part of the semantics. The reference companion goes further
in restriction: it has no equality transport and refuses all stale epochs or
cross-anchor reuse. Its tests validate that refusal policy, not a completed
implementation of the more general transport design.

\subsection{Different failures must remain distinguishable}
Heather should distinguish at least four broad outcomes. A demand can have an
accepted proof; it can remain unresolved within a profile; a provider or
transport can be refused because its schema or context does not match; or the
negation of the actual demand can be proved. These outcomes must not share a
single undifferentiated failure label.

For behavioral requests, an abstract counterexample is an additional intermediate
status. Its promotion to concrete refutation requires the bridge developed
below. Budget exhaustion, an unsupported grammar, a corrupted certificate,
and an unrealizable abstract assignment are not interchangeable forms of
``the program is wrong''.

\section{Realizable affine frames for behavioral synthesis}\label{sec:frames}

\subsection{The gap in unrestricted coefficient comparison}
The round-three reports develop affine models of list lengths and identify
two hazards: a positive abstract length may not be realizable for an empty
element type, and independently possible observations may not be jointly
realizable under a relational input guard \cite{round3a,round3b}.

A useful response is not to abandon finite reasoning. It is to attach the
finite reasoning to the \emph{actual input space}. For an admissibility
predicate $D:A\to\mathsf{Prop}$, define
\[
 S=\{x:A\mid D(x)\},\qquad o:S\to\Q^k.
\]
The type $S$ includes the full input configuration and its guard. It may encode
multiple correlated inputs, an empty element type, an invariant, or a chosen
provider environment. The observation $o$ is interpreted jointly on that type.

\subsection{Definition and determining theorem}
\begin{definition}[Realizable affine frame]
For a nonempty admissible input type $S$ and observation $o:S\to\Q^k$, a
realizable affine frame consists of points $s_0,s_1,\ldots,s_r\in S$ and a
coverage theorem
\begin{equation}\label{eq:frame-cover}
 \forall s\in S,\quad
 o(s)-o(s_0)\in
 \operatorname{span}_{\Q}\{o(s_j)-o(s_0):1\le j\le r\}.
\end{equation}
The frame need not be linearly independent. Its points are actual admissible
inputs, not merely satisfying assignments in an over-approximate abstract model.
\end{definition}

\begin{theorem}[Domain-sensitive affine determination]\label{thm:frame}
Let $(s_0,\ldots,s_r)$ be a realizable affine frame for $o:S\to\Q^k$.
For any affine maps $F,G:\Q^k\to\Q$, the following are equivalent:
\begin{align}
 &\forall s\in S,\ F(o(s))=G(o(s)),\label{eq:all-inputs}\\
 &\forall j\in\{0,\ldots,r\},\ F(o(s_j))=G(o(s_j)).\label{eq:frame-inputs}
\end{align}
The same assertion holds for affine maps to a vector space over $\Q$.
\end{theorem}
\begin{proof}
The forward direction follows because each $s_j$ belongs to $S$. For the
reverse direction, set $H=F-G$ and write $H(v)=c+L(v)$, where $L$ is linear.
Equation~\eqref{eq:frame-inputs} implies $H(o(s_0))=0$ and
\[
 L(o(s_j)-o(s_0))=H(o(s_j))-H(o(s_0))=0.
\]
For any $s$, coverage expresses $o(s)-o(s_0)$ as a linear combination of these
differences. Linearity gives $L(o(s)-o(s_0))=0$. Therefore
\[
 H(o(s))=H(o(s_0))+L(o(s)-o(s_0))=0,
\]
as required. The argument only uses addition and rational scalar multiplication
in the codomain, so it also applies to vector-valued affine maps.
\end{proof}

The two directions explain exactly where realizability enters. Coverage of the
source image is sufficient for positive transfer from a spanning family of
abstract points. Requiring that the family itself is realized additionally
makes failure at one of those points a source-level counterexample, and yields
the equivalence above. A certificate should not conflate the two requirements.

\subsection{Lifting from a model to the exact program}
Let $e$ be a candidate in a specified expression grammar, with concrete
semantics $\ev(e,s)$ and observation $\ell$ of its result. Suppose a proved
model theorem supplies an affine map $F_e$ satisfying
\begin{equation}\label{eq:model-bridge}
 \forall s\in S,\quad \ell(\ev(e,s))=F_e(o(s)).
\end{equation}
Suppose a requested target observation is represented by $G$. Then
Theorem~\ref{thm:frame} reduces the universal contract
\[
 \forall s\in S,\quad\ell(\ev(e,s))=G(o(s))
\]
to the frame equalities. A failed frame equality gives the concrete admissible
input $s_j$, not merely an abstract tuple.

To accept a rendered Lean candidate $f$, one more theorem is needed:
\[
 \forall s\in S,\quad f(s)=\ev(e,s),
\]
or an adequate direct equality of the requested observations. The expression
that was modeled must be the expression the user actually selected. A proof
about a convenient surrogate, or a provider's claimed length law, does not
close this last bridge.

\subsection{Four exact list-domain instances}
For lists, observe their lengths as nonnegative integers embedded in $\Q$.
The following frames have elementary coverage and realization proofs.

\begin{center}\small
\begin{tabularx}{\textwidth}{@{}P{0.24\textwidth}P{0.27\textwidth}Y@{}}
\toprule Admissible inputs & Observed frame & Why it covers and is realized\\\midrule
$k$ independent lists over an inhabited type & $0,e_1,\ldots,e_k$ & Every length
tuple is a sum of basis directions; empty lists and singleton lists realize the
points using a supplied element.\\
$k$ lists over an empty type & $0$ only & Every list is empty, so the image is a
singleton. No positive length is admissible.\\
Two inhabited-type lists with equal lengths & $(0,0),(1,1)$ & Every observation
is $n(1,1)$; paired empty and paired singleton lists realize the frame.\\
$k$ inhabited-type lists of even lengths & $0,2e_1,\ldots,2e_k$ & Every admitted
tuple is a nonnegative integer combination of these directions; length-two
lists realize them.\\
\bottomrule
\end{tabularx}
\end{center}

For example, the length models of
\[
 (x,y)\longmapsto x\mathbin{+\!+}y
 \quad\text{and}\quad
 (x,y)\longmapsto y\mathbin{+\!+}y
\]
are $n_1+n_2$ and $2n_2$. They are unequal on all of $\N^2$ but equal on the
diagonal domain $n_1=n_2$. Testing $(1,0)$ as a negative input in that domain
would discard the admissibility guard. The frame instead uses $(1,1)$, a pair
of actual lists satisfying the guard.

Likewise, constant-empty output and the identity have models $0$ and $n$.
On $\mathsf{List}(\mathsf{Empty})$ their source behavior is identical and the
one-point frame proves the observation law. On $\mathsf{List}(\N)$ the singleton
input is a concrete negative witness. This is a change in the domain of the
specification, not a change in the soundness of arithmetic.

An empty \emph{admissible input type} $S$ is a separate case from an empty
\emph{element type}. When $S$ is empty, every universal contract is vacuous,
proved from an emptiness theorem, and there is no base point $s_0$. By contrast,
$\mathsf{List}(\mathsf{Empty})$ contains the empty list and has a valid
one-point frame. A native interface should have distinct constructors for
these cases.

\subsection{Why affine length models are valid for the chosen grammar}
Consider the list-expression grammar
\[
 e::=x_i\mid []\mid [a]\mid e_1\mathbin{+\!+}e_2
       \mid\operatorname{reverse}(e),
\]
where the singleton constructor is available only when its element $a$ is
supplied. Assign the models
\[
 x_i\mapsto n_i,\quad []\mapsto0,\quad [a]\mapsto1,\quad
 e_1\mathbin{+\!+}e_2\mapsto F_{e_1}+F_{e_2},\quad
 \operatorname{reverse}(e)\mapsto F_e.
\]

\begin{lemma}[Length interpretation]\label{lem:length}
Every expression in this grammar has an affine length model
$F_e(n)=c_e+\sum_i a_{e,i}n_i$ with natural coefficients, and
$|\ev(e,\rho)|=F_e(|\rho_1|,\ldots,|\rho_k|)$ for every well-typed input
environment $\rho$.
\end{lemma}
\begin{proof}
Induct on $e$. An input has its observed length. Empty and singleton lists have
lengths zero and one. The length of an append is the sum of lengths, so adding
the induction hypotheses adds constants and coefficients. Reversal preserves
length, so it preserves the induction hypothesis unchanged. Each constructor
preserves the asserted affine form and its interpretation.
\end{proof}

Differences of two natural-coefficient models are interpreted over $\Q$ when
applying Theorem~\ref{thm:frame}. Equality of their rationally embedded lengths
is equivalent to equality of the original natural lengths. This scalar change
is an explicit model bridge, not an implicit reinterpretation of the source
program's arithmetic.

\subsection{The number of observations is controlled by rank}
\begin{proposition}[Size and minimality of a determining family]\label{prop:rank}
Let $S$ be nonempty and let the affine hull of $o(S)\subseteq\Q^k$ have
dimension $r$. There exists a realizable affine frame of $r+1$ points.
No family with fewer than $r+1$ observed points determines the restrictions of
\emph{all} affine maps $\Q^k\to\Q$ to $o(S)$ by equality of their values.
\end{proposition}
\begin{proof}
Choose $s_0\in S$. The vector space spanned by the differences
$o(s)-o(s_0)$ has dimension $r$. A basis can be selected from this spanning
set, giving $r$ further realized points and the required coverage.

Now take $m<r+1$ proposed observed points. If $m=0$, two different constant maps
already show failure. Otherwise their affine hull has dimension at most
$m-1<r$. It is a proper affine subspace of the affine hull of $o(S)$.
Choose an affine functional that vanishes on the smaller hull but not on the
larger one: translate by one selected point, extend a basis of the smaller
difference space, and take a coordinate functional on an added basis vector.
Extend it to $\Q^k$. This functional and zero agree on all selected points but
disagree at some point of $o(S)$, proving the lower bound.
\end{proof}

This is an existence and information-content statement, not an algorithm for
deciding arbitrary admissibility predicates or discovering their affine hulls.
The proof can use mathematical choices that a computational interface must
supply separately. Nor is the lower bound necessarily sharp for a smaller
candidate grammar: it concerns all affine maps, not only the maps a particular
synthesizer happens to produce.

Nevertheless, the proposition suggests a useful engineering principle. Input
guards can lower the dimension of the observation space and hence reduce both
the certificate size and the number of relevant negative cases. Heather should
model that reduced space rather than independently abstracting each input and
then trying to repair impossible counterexamples afterward.

\subsection{What the theorem does not authorize}
Affine determination decides \emph{equalities}. Agreement with an inequality
at a base point and basis points is not a universal nonnegativity proof: the
affine function $1-n$ is nonnegative at $n=0,1$ but negative at $n=2$.
A positivity service needs a different certificate, for example cone-based
reasoning with its own coverage and sign conditions.

Similarly, length equality proves neither permutation nor sortedness. A
function that branches on an unmodeled property of its input is outside the
chosen grammar. Finite observations of that unrestricted function cannot be
promoted by borrowing the affine theorem for another grammar. A future
extension can use other finite-dimensional feature models, but each requires
its own exact denotation and coverage proofs.

\section{Leant integration: from candidates to retained behavior}\label{sec:leant}

\subsection{A request contract}
Heather should submit a construction request only after the intended type and
behavioral specification have been fixed. A request includes the elaborated
candidate type, the proposition about the candidate, selected structures,
local assumptions, admissible input domain, permitted providers, resource
limits, and the original source location. For a length request, it additionally
identifies the observation grammar and any domain-frame package to be used.

The result must distinguish a candidate that has merely passed a callback from
a candidate with retained behavioral evidence. Leant already makes the former
boundary explicit in its \code{Verified} abstraction \cite{leant-verification}.
Heather's stronger acceptance object is conceptually
\[
 (f,p),\qquad f:T,\quad p:\Spec(f),
\]
plus the environmental and source-association information needed to replay
those terms. A host-language receipt may refer to the Lean declaration or proof
artifact, but the receipt's constructor is not a replacement for that proof.

The nominal epochs and exact-occurrence handling in Leant's post-verification
stage are useful foundations for routing this record. They should be retained,
not flattened to a string-to-Boolean map. Association checks and proof checks
remain separate obligations \cite{leant-post}.

\subsection{A concrete affine handoff}
For the restricted list fragment, the proposed pipeline is as follows. Parse
and elaborate the selected candidate once. Produce a typed grammar expression
$e$ and a correspondence theorem for that exact candidate. Derive its length
model using the grammar's interpretation theorem. Select a proved frame for
the actual admissible input type. Check the finite affine equalities or
reconstruct an admissible source counterexample. Finally, retain the proof of
the original specification, or the proof of its negation, at the selected
candidate.

The companion implements the parsing, model computation, four fixed frames,
certificate reconstruction, and concrete negative evaluation in Python. It
emits Lean source candidates for positive cases. It does not implement the
Lean correspondence theorem or connect its model to Leant's Exference term
graph. That missing correspondence is an explicit integration gate, not a
minor matter of converting a JSON response.

The current Length adapter already distinguishes source association from
model-relative behavior and treats product-domain queries separately
\cite{leant-adapter}. Heather's realizable-frame package can sit after that
boundary, provided its model is proved to describe the same source terms and
provider environment. It should not bypass the existing handoff's refusal
conditions.

\subsection{Provider laws must be attached to their providers}
Suppose a provider $h$ is abstracted by a length law
$|h(x)|=a|x|+b$. A passive record containing $(a,b)$ is a hypothesis about the
model, not evidence that this exact $h$ satisfies the law. Heather needs a Lean
theorem
\[
 \forall x,\quad |h(x)|=a|x|+b
\]
under the admitted domain and structures, together with the mapping from the
provider identifier to that $h$.

There are two honest alternatives when this proof is unavailable. Search may
use the law heuristically but may not promote its conclusion to a theorem; or
the published result may remain conditional on the provider law. The latter
is legitimate mathematics when the assumption is visible. It is not an
unconditional verification of the provider.

\subsection{Negative outcomes and unfinished sketches}
A negative result can be an unsupported model, an exhausted search, a model
counterexample, a realized source counterexample, or a proof of the negation
of the original specification. Only the last two, connected to the exact
candidate, justify source-level refutation. Even a perfectly reconstructed
abstract solver assignment remains model-relative when a provider law has not
been established.

For a sketch with possible completions $C$, one failed completion $c_0$ proves
only $\neg\Spec(c_0)$. Rejecting the entire sketch requires a proof such as
\[
 \forall c\in C,\quad\neg\Spec(c),
\]
not merely a counterexample to one candidate chosen by the synthesizer.
Backward requirement propagation may rule out a class of completions, but the
coverage of that class is itself part of the proof.

\subsection{Replay the actual edit}
The final editor action must replay the exact selected source replacement in
the destination context and verify all resulting proof obligations. Rechecking
a semantically intended candidate while inserting a different rendered variant
is insufficient. The retained record identifies the interpreted statement,
selected candidate, exported proof, dependencies, and policy at acceptance.
On a toolchain or import change, replay creates a new record rather than
relabeling an old one as current.

\section{Certified computer algebra as an interface service}\label{sec:cas}

\subsection{Three acceptable routes and one forbidden shortcut}
Heather can use an algorithm with a proved correctness theorem, an untrusted
algorithm that supplies a certificate to a proved checker, or a tactic that
constructs an ordinary Lean proof. The acceptance boundary in every case is
the requested relation on the original objects. A numeric answer, a solver's
success flag, or an equality about a different reified expression is not enough.

For a certificate route, the shape is
\[
 \operatorname{check}(x,y,c)=\mathsf{true}
 \quad\Longrightarrow\quad R(x,y).
\]
The source-to-expression bridge and the proposition that the actual checker
accepted the actual certificate are additional obligations. A checker soundness
theorem alone says nothing about an untrusted native process's returned bytes.
The default route should use kernel reduction where practical or reconstruct
an ordinary proof. Accelerated native evaluation needs its own explicit,
versioned trust policy and axiom audit \cite{round3a,round3b,lean-axioms}.

The forbidden shortcut is to place arbitrary computer-algebra simplification
inside definitional equality. A false or incomplete result must fail as a proof
attempt, not change what the kernel means by type conversion.

\subsection{Typed reification rather than one untyped universal AST}
A reusable arithmetic reifier should return a syntax tree together with its
interpretation bridge in a fixed algebraic structure. The same ring-expression
syntax can support several compatible checkers. It should not erase the
distinction between natural-number length arithmetic, integer exact division,
field inversion, and real inequalities merely because all four print as
familiar arithmetic.

For example, a polynomial identity checker can validate a certificate
$p=\sum_i c_iq_i$ by exact normalization under a fixed commutative-ring
interpretation. Transport through a ring homomorphism then uses its theorem.
The reverse implication for ideal membership need not hold after extending
scalars. Likewise, casting an integer quotient is not automatically field
division; Section~\ref{sec:frey} isolates the missing interface.

The first implementation should reuse existing Lean arithmetic tools before
building a larger checker. The round-three Schist checker is a useful
historical seed for a tiny affine equality chain, not evidence that a
multivariate polynomial reifier is already available \cite{round3a,round3b}.
The practical question is whether an external service supplies capabilities or
reusable certificates that equally equipped Lean does not already provide.

\subsection{Result strength and quantitative indices}
A service request fixes the kind of result. A witness of $P(x)$ is different
from a complete characterization of all $x$ satisfying $P$. A factorization
identity is different from a factorization together with irreducibility and
normalization guarantees. An upper error bound is different from a sharp
constant. Heather represents these distinctions by predicates and dependent
indices, not a single ladder of result confidence.

Finite precision can also be part of an interface. For formal series, define
$f\equiv_Ng$ to mean equality of coefficients of degrees below $N$. For $N\ge1$,
formal differentiation gives
\[
 f\equiv_Ng\quad\Longrightarrow\quad f'\equiv_{N-1}g'.
\]
Indeed, coefficient $j$ of a derivative is $(j+1)$ times coefficient $j+1$ of
the original series. This is a theorem-backed precision change, not an implicit
loss of information hidden by the notation. Integration needs a separate
interface with the necessary scalar-division assumptions.

Similarly, if $f$ is $L$-Lipschitz, $d(x,x_0)\le\varepsilon$, and an approximation
$y_0$ satisfies $d(f(x_0),y_0)\le\delta$, the triangle inequality gives
\[
 d(f(x),y_0)\le L\varepsilon+\delta.
\]
The source can request this enclosure; Heather propagates the quantitative
requirements and retains the bound proof. It may not render the enclosure as
an exact equality. These examples show how richer type interfaces can carry
behavioral constraints without new kernel-level logic.

\section{Worked mathematics I: exact quotients without vacuity}\label{sec:frey}

\subsection{The inherited benchmark}
The two syntheses recommend exact quotient transport as a small, non-vacuous
arithmetic use site. Their checked helper development is historical evidence
reported by those documents, not a fresh compilation in this archive
\cite{round3a,round3b}. The mathematical problem is the following.

Let $a,b\in\Z$ and $p\in\N$ satisfy
\begin{equation}\label{eq:frey-assumptions}
 p\ge4,\qquad p\text{ odd},\qquad a\equiv3\pmod4,\qquad 2\mid b.
\end{equation}
Set
\[
 N_2=b^p-1-a^p,\qquad N_4=-a^pb^p,
 \qquad q_2=N_2/4,\quad q_4=N_4/16,
\]
where the quotients are initially integer division. The goals are exactness
of those divisions and the corresponding rational cast identities.

\begin{proposition}[Operation-specific arithmetic premises]
Under \eqref{eq:frey-assumptions}, $4\mid N_2$ and $16\mid N_4$. The second
claim requires only $p\ge4$ and $2\mid b$. Consequently the integer quotients
satisfy $4q_2=N_2$ and $16q_4=N_4$ and become the corresponding quotients in
$\Q$ after casting.
\end{proposition}
\begin{proof}
Write $b=2t$. Since $p\ge4$, $b^p=2^pt^p$ is divisible by $16$ and hence by
$4$. Since $a\equiv-1\pmod4$ and $p$ is odd, $a^p\equiv-1\pmod4$.
Thus $N_2\equiv0-1-(-1)=0\pmod4$. Also $16\mid b^p$ implies
$16\mid-a^pb^p=N_4$, independently of the remaining assumptions.
Divisibility makes the integer quotients exact. A ring homomorphism from
$\Z$ to $\Q$ preserves the two balance equations. Since $4$ and $16$ are
invertible in $\Q$, solving those equations gives the claimed rational
quotient identities.
\end{proof}

\subsection{The type-level interface}
An exact-quotient package should expose a balance theorem rather than a vague
``division is safe'' flag. One possible result interface is
\[
 \mathsf{ExactQuot}(n,d)=\{q:\Z\mid d q=n\}.
\]
A construction from integer division carries its divisibility proof. Casting
the balance equation along $\iota:\Z\to K$ produces
$\iota(d)\iota(q)=\iota(n)$. Dividing in $K$ requires a unit witness for
$\iota(d)$, or a theorem supplying it from the destination structure.

Nonzeroness in $\Z$, regularity in an arbitrary ring, and invertibility in the
destination are different facts. A field of characteristic zero is one
sufficient destination package; it is not the weakest possible condition.
The actual consumer asks for the unit of the denominator's image.

A proposed Heather trace is:
\par\noindent\begin{minipage}{\linewidth}
\begin{lstlisting}
establish q2_exact : 4 * q2 = N2
  via integer exact division
establish q2_cast : cast(q2) = cast(N2) / 4
  via transport of q2_exact to Rational
\end{lstlisting}
\end{minipage}\par
The first method asks for $4\mid N_2$ and the relevant division interface; the
second asks for the fixed cast map and the destination unit. The proof engine
may derive those premises from the selected arithmetic helper theorems.
It may not replace the integer quotient expression by a rational one before
justifying their relationship.

\subsection{Satisfiable support and diagnostic mutations}
The tuple $(a,b,p)=(3,2,5)$ satisfies the assumptions and gives
\[
 N_2=-212,\quad q_2=-53,\qquad N_4=-7776,\quad q_4=-486.
\]
This witness matters because a benchmark run only inside an inconsistent
Fermat counterexample package could discharge every goal by contradiction,
without exercising any exact-quotient reasoning \cite{round3a,round3b}.

Removing oddness permits $(3,2,4)$, where $N_2=-66$ is not divisible by $4$.
Removing the evenness condition permits $(3,3,5)$, where $N_2=-1$.
Weakening the exponent bound permits $(3,2,3)$, where $N_4=-216$ is not divisible
by $16$. Removing the residue condition permits $(1,2,5)$, where $N_2=30$.
These are concrete mathematical counterexamples to the corresponding weakened
claims, not merely failed tactic invocations.

For a required arithmetic route, Heather should elaborate its helper theorem
in a clean, dependency-closed telescope containing the advertised arithmetic
assumptions. This prevents accidental use of unrelated local contradictions.
A syntactic dependency audit can enforce that support restriction; it does not
establish that every allowed premise is mathematically indispensable, nor
that a theorem constant appearing decoratively constitutes the advertised
method.

The executed Nat use-site demonstration in this archive is deliberately
smaller: it inserts a nonzero proof from positivity and retrieves a divisibility
assumption. It exports a proof of the two readiness premises. It neither
constructs the Frey quotients nor proves the integer-to-rational casts.

\section{Worked mathematics II: locality and analysis}\label{sec:analysis}

\subsection{The autonomous-derivative theorem in full}
Let $U\subseteq\R$ be open, let $W,\varphi:\R\to\R$, and let
$G_n,H_n:\R\to\R$ be families. Assume that for every $x\in U$,
$W$ is differentiable at $x$ with derivative $\varphi(W(x))$; for every $n$
and $x\in U$, $G_n$ is differentiable at $W(x)$ with derivative $H_n(W(x))$;
and
\[
 G_0(W(x))=W(x),\qquad
 G_{n+1}(W(x))=H_n(W(x))\varphi(W(x)).
\]
These are the mathematical roles present in the ProveIt theorem, with the
supplied derivative family renamed $H$ to avoid confusing a name with a
derivative operator \cite{autonomous}.

\begin{theorem}[Range-local autonomous differentiation]
For all $n\in\N$ and $x\in U$,
\[
 D^nW(x)=G_n(W(x)).
\]
\end{theorem}
\begin{proof}
Induct on $n$, retaining the universal quantifier over $x\in U$. For $n=0$,
$D^0W(x)=W(x)=G_0(W(x))$ by the initial equation.

Assume the assertion for $n$. Fix $x\in U$. Since $U$ is open, it contains a
neighborhood of $x$. On that neighborhood the induction hypothesis identifies
$D^nW$ with $G_n\circ W$. The latter is differentiable at $x$ by the chain
rule, with derivative
\[
 H_n(W(x))\varphi(W(x))=G_{n+1}(W(x)).
\]
Equality on a neighborhood transfers this derivative to $D^nW$. Its derivative
at $x$ is $D^{n+1}W(x)$, proving the successor assertion.
\end{proof}

The proof does not require $G_n$ to be differentiable away from $W(U)$.
It does not differentiate a merely pointwise equality at one point. In a Lean
translation, differentiability evidence must be retained rather than inferring
it from equality of totalized derivative values alone.

\subsection{The consumer interface and its failure neighbor}
The critical method is a local derivative transfer. Its inputs are two fixed
functions, a point, a neighborhood-equality proof, and an appropriate derivative
statement. A convenience adapter obtains neighborhood equality from equality
on an open set plus membership. The source file already performs exactly this
composition using existing Lean lemmas \cite{autonomous}.

Heather's contribution is to infer those adapters under the operation's
requirements and explain failures. If openness is removed, it should report
that the selected route lacks a neighborhood; another route may still establish
one. If only $f(x)=g(x)$ is supplied, differentiation is invalid: $f(t)=t$ and
$g(t)=0$ agree at zero but have different derivatives there. If the induction
hypothesis was prematurely specialized to the selected point, Heather should
identify that missing uniformity rather than ask the author to experiment
with arbitrary tactic variants.

\subsection{Almost-everywhere equality has different consumers}
The corresponding integration adapter has a different relation. Mathlib's
Bochner integral API provides an almost-everywhere congruence theorem for a
fixed measure. Its total integral definition and congruence interface should
not be confused with separate integrability requirements needed by other
operations \cite{bochner}.

Thus an evidence row containing $f=g$ almost everywhere with respect to $\mu$
can discharge an integral-congruence requirement for $\mu$, but it cannot
justify evaluation at a selected point or an integral with respect to an
unrelated measure. Equality off a singleton illustrates the distinction:
under Lebesgue measure a point can be negligible, whereas under a Dirac measure
at that point it is not. The measure is part of the proposition's identity.

Likewise, exchanging a limit with an integral requires the premises of an
applicable interchange theorem. Domination is one sufficient route, not a
universal necessary condition for all possible interchange arguments. Heather
should report the premises of the chosen route without overgeneralizing the
reason for a particular refusal.

\section{Worked mathematics III: a whole construction carries laws}\label{sec:tensor}

For a finite family $(H_i)_{i\in I}$ of commutative $A$-algebras, consider the
functor
\[
 T\longmapsto\prod_{i\in I}\operatorname{Hom}_{A\text{-alg}}(H_i,T).
\]
A finite tensor construction represents this functor. In the cited Fermat
source, each $H_i$ is additionally finite and free as an $A$-module, and the
constructed representing algebra has those properties \cite{tensor}.

The proof's essential recursive interface is precise. The empty family is
represented by $A$. If $W_0$ represents a smaller family, then
$W_0\otimes_A H$ represents that family with one additional factor. Given two
algebra maps into a commutative target, the tensor universal property provides
the combined map. The equivalence of homomorphism spaces is natural because
postcomposition commutes with restriction to each factor. Finite freeness is
preserved at the tensor step: with finite bases of sizes $m$ and $n$, the pure
tensors of basis vectors give a basis of size $mn$. This supplies a mathematical
route for the relevant closure property, while the source uses Lean's module
interfaces to assemble it \cite{tensor}.

A useful Heather package would expose the representing equivalence,
naturality, and module properties as a dependent interface for the one
constructed $W$. It would not independently synthesize one $W$ for finiteness,
another for representability, and then merge their rows by printed name.
Its naturality facet would quantify over target algebras, their structures,
postcomposition maps, and the actual representing equivalence.

The missing-factor hypothesis has a simple negative neighbor. For a singleton
family over $A=\Q$ with $H=\Q[X]$, the functor is already represented by $H$.
A natural representing isomorphism forces any other representing algebra to
be isomorphic to $H$: apply the natural isomorphism and its inverse to identity
maps to obtain mutually inverse algebra maps. But the powers
$1,X,X^2,\ldots$ are linearly independent, so $H$ is not a finite-dimensional
$\Q$-module. Finiteness of the index set alone therefore cannot justify the
finite-module facet.

This example also separates existential mathematical construction from
executable data extraction. The source theorem is an existence statement.
A proof may introduce its existential witness within a proposition-valued
argument. An API promising an executable enumeration, explicit bases, or a
computable normalization routine needs additional data, or an explicitly
noncomputable construction policy. A proposition asserting finiteness is not
itself a program that enumerates the carrier.

\section{Relation to existing type-system ideas}

Refinement types and type-directed proof obligations are established ideas,
not Heather inventions. Liquid Types combine type inference with predicate
abstraction to infer useful refinements automatically. Heather shares the
interest in finite, disciplined inference, but it does not claim to make
arbitrary Lean propositions decidable or to infer a principal refinement for
a mathematical development. Its automation is deliberately incomplete outside
a selected, theorem-backed fragment \cite{liquid}.

Occurrence typing makes the classification of an expression depend on the
facts available at its occurrence, particularly in control flow. Heather's
branch-local facets have a related motivation: a hypothesis learned in one
branch changes what can be established there, not what may be assumed in its
sibling. Heather uses retained Lean evidence and explicit dependent contexts
rather than treating this analogy as a ready-made implementation
\cite{occurrence}.

Within Lean, automatic parameters, property tactics, dependent structures, and
program-verification tools already cover important pieces of the proposed
interface. Heather's research question is whether a uniform, mathematically
legible protocol for those pieces improves proof construction and repair.
The strongest alternative is not bare tactic syntax with no libraries; it is
well-engineered ordinary Lean with the same semantic packages and services
\cite{lean-application,funprop,mvcgen}.

The two round-three reports make this distinction central. Their consensus
calculus is a foundation to reuse; their implementation questions are where
Heather seeks to add specificity. The realizable-frame package, in particular,
is a domain-adaptive use of elementary finite-dimensional linear algebra, not
a claim that finite testing proves unrestricted programs correct
\cite{round3a,round3b}.

\section{How far should Lean's type system be extended?}\label{sec:kernel}

\subsection{Four different meanings of ``extend''}
Heather distinguishes four interventions that are often conflated. A library
extension adds predicates, bundled structures, and theorems. An elaboration
extension changes how surface types and omitted arguments are interpreted.
A kernel extension adds new primitive term or type constructors with checking
rules. A conversion extension changes which expressions count as
definitionally equal.

The first two already suffice to express and implement the proposed mathematical
interfaces. Lean's existing dependent types can state the contracts; Heather
changes which evidence is routinely carried to a use site and how it is shown.
That is a real extension of the author's typing discipline, but not an increase
in the logic's expressive power \cite{lean-reference,round3a,round3b}.

\subsection{A precise primitive refinement alternative}
A plausible kernel experiment would add, for $A:\mathsf{Type}_u$ and
$P:A\to\mathsf{Prop}$, a primitive $\mathsf{Ref}(A,P)$ with operations
\[
 \frac{a:A\quad h:P(a)}{\mathsf{pack}(a,h):\mathsf{Ref}(A,P)},
\]
\[
 r:\mathsf{Ref}(A,P)\Longrightarrow
 \mathsf{value}(r):A,\qquad
 \mathsf{evidence}(r):P(\mathsf{value}(r)),
\]
and reduction rules
\[
 \mathsf{value}(\mathsf{pack}(a,h))\equiv a,
 \qquad
 \mathsf{evidence}(\mathsf{pack}(a,h))\equiv h.
\]
Subsumption would still require an explicit proof. Given
$k:\forall x,\ P(x)\to Q(x)$, define
\[
 \mathsf{weaken}_k(r)=
 \mathsf{pack}(\mathsf{value}(r),
               k(\mathsf{value}(r),\mathsf{evidence}(r))).
\]
This keeps the underlying value fixed. It does not insert an oracle for
$P\Rightarrow Q$ into type checking.

The evident translation is to Lean's existing subtype
$\{a:A\mid P(a)\}$, with its ordinary constructor and projections. On these
rules the primitive is therefore a candidate representation optimization, not
new mathematics. Introduction still requires $h$: making the evidence
unnecessary would either move the proof obligation elsewhere or change the
logic. The implementation experiment would need to establish preservation of
typing and reduction under translation, alongside the ordinary metatheoretic
and trusted-code obligations of a kernel change.

No measured encoding obstruction in the inspected examples demands this
primitive. Local facets avoid repeated packaging already. Before implementing
it, profile ordinary subtypes, structures, proof sharing, reducible adapters,
and existing elaborator facilities on the same workload. A native constructor
must demonstrate a benefit those alternatives cannot provide.

\subsection{Why not semantic subtyping in conversion?}
A rule that silently accepts $A\mid P$ wherever $A\mid Q$ is expected whenever
$\forall x,\ P(x)\to Q(x)$ is true would couple conversion to an unrestricted
mathematical entailment problem. A bounded prover can implement a useful
\emph{elaboration attempt}, but its timeout cannot decide that the subtype
relation is false. A kernel whose conversion depends on such attempts would
make a basic checking operation sensitive to search policy.

Equality reflection creates a related problem: a propositional equality would
be promoted into conversion rather than used by an explicit transport term.
Heather needs proof-directed transport, not that promotion. Keeping the proof
step explicit in the core preserves a small acceptance boundary while allowing
its routine surface syntax to be omitted.

Restricted new definitional laws, such as laws intended to reduce transport or
functoriality overhead, remain a research option. They must be justified by a
specific cost and a metatheory for the exact rule set. Heather does not infer
that they are harmless merely because their propositional counterparts are
provable.

\subsection{Behavioral and graded types without a new conversion theory}
A function can be exported as a subtype of functions satisfying a contract,
a structure with a function field and laws, or an unchanged function with
separate evidence parameters. The first two are useful at module boundaries;
the last is preferable for accumulating local knowledge. These are different
interfaces for the same logical content, not literally the same Lean term
encoding.

Quantitative guarantees likewise fit dependent indices: a finite jet's order,
an enclosure's radius, a Lipschitz constant, or a permitted domain. Operations
have theorems transforming these indices. Their transformations need not all
form a single semiring or scalar strength order; heterogeneous mathematical
relations deserve heterogeneous indices.

For resource-sensitive programs, state-transition specifications or an
appropriate linear interface may be useful. That does not make mathematical
proofs themselves single-use resources. A host receipt's epoch can regulate
association with one request, while a theorem remains duplicable evidence.
Heather keeps these two notions of constraint separate.

\section{The delivered executable experiment}\label{sec:experiment}

\subsection{What exists}
The archive contains a standard-library-only Python implementation and five
small source examples. Its implemented language is intentionally smaller than
the proposed narrative surface:
\par\noindent\begin{minipage}{\linewidth}
\begin{lstlisting}
heather diagonal_balance
domain diagonal
inputs xs, ys
claim length(append(xs, ys)) == length(append(ys, ys))
\end{lstlisting}
\end{minipage}\par
It parses the source into a bounded expression tree, checks the available list
constructors and inputs, computes affine length models, selects a fixed
domain-specific frame, constructs a result certificate, and rechecks that
certificate against the original request. Positive results emit ordinary Lean
proof candidates; negative results include actual lists whose evaluations have
different lengths.

A separate small context model resolves the two premises of an exact-quotient
consumer over $\N$. It uses only positivity/nonzeroness implications and
explicit divisibility assumptions. This is enough to exercise proof-argument
insertion, exact object matching, absent-premise behavior, and snapshot
refusal. It is not a general rule registry or exact-division implementation.

\subsection{What ran}
The delivered run used Python 3.13.5 and passed all 18 test methods. The log and
machine-readable results are included. Counts are separated by what they
measure:
\begin{center}\small
\begin{tabularx}{\textwidth}{@{}Y r@{}}
\toprule Executed check category & Count\\\midrule
Test methods completed with no failures or errors & 18\\
Concrete expression-length versus affine-model comparisons & 7,500\\
Admissible-domain evaluations in generated law checks & 7,200\\
Concrete negative witnesses checked in generated examples & 466\\
Deliberately corrupted result certificates refused & 14\\
New Lean compilations performed & 0\\
\bottomrule
\end{tabularx}
\end{center}
The generated cases use fixed random seeds and small finite length ranges.
They are not 14,700 proved theorems. The universal mathematical argument for
the intended algorithm is Lemma~\ref{lem:length} with
Theorem~\ref{thm:frame}; the correspondence between the Python program and that
argument has not been mechanically verified.

The five named source examples produce four positive model conclusions and
one concrete counterexample. Constant-empty output preserves length over
\code{List Empty}, but not over \code{List Nat}. The diagonal contract is
accepted under equal input lengths; removing the guard produces a concrete
counterexample. Reversal/append and even-domain reversal examples exercise
ordinary positive cases. Constructing a singleton of \code{Empty} is refused
as a typing error, not mislabeled as a mathematical counterexample.

The use-site demonstration resolves two requirements with one implication-rule
firing and three retained evidence entries, counting the two assumptions. These
are trace statistics for a toy example. They are not measurements of a native
Lean evidence index, proof-term growth on a real file, or author productivity.

\subsection{Trust and implementation boundaries}
The certificate rechecker reconstructs model and frame information instead of
trusting serialized coefficients or a result flag. It also evaluates the
actual expression on a negative witness. However, producer and checker share
parser, expression, and model code. They are not independent implementations,
and neither is a verified Lean checker.

The generated Lean files are labeled uncompiled. A working Lean toolchain was
not available in this environment, and an attempted toolchain retrieval did
not succeed. The archive therefore makes no new kernel-checking claim. It
contains reproducible commands for attempting compilation in an appropriate
Mathlib project rather than a fabricated successful build log.

General expression rewriting, checked context migration, automatic frame
discovery, arbitrary provider contracts, a full \code{requires} parser, and the
Leant handoff remain design specifications. The Python context uses exact
anchors and refuses transport. That narrower behavior is a deliberate boundary
of the experiment, not evidence that the more general operations have been
implemented.

\subsection{Why this still tests a useful design question}
The experiment establishes an executable counterpoint to coefficient equality
on an unrestricted observation space. The same source grammar gives different
correct outcomes when its domain changes; impossible abstract witnesses do
not become source refutations; and corrupted or mismatched certificates are
not accepted merely because a producer labeled them successful.

These results justify attempting a native semantic bridge for the small
fragment. They do not yet justify a new proof language. The next gate is a
Lean-checked interpretation theorem and exact-candidate correspondence,
followed by a genuine comparison against the strongest ordinary-Lean route.

\section{Adversarial acceptance and the evaluation contract}\label{sec:evaluation}

\subsection{Positive neighbors prevent a trivial refusal system}
Every negative test needs a nearby admitted case and an expected failure kind.
Otherwise an implementation that rejects everything can appear safe. The
following table separates what the companion executes from important native
tests still needed.
\begin{center}\small
\begin{tabularx}{\textwidth}{@{}P{0.29\textwidth}Y P{0.20\textwidth}@{}}
\toprule Mutation & Correct response and positive neighbor & Status\\\midrule
Remove divisibility from a use site & Leave that premise unresolved; resolve it
when the matching assumption is supplied. & Executed toy resolver.\\
Reuse positivity for another object & Refuse cross-anchor reuse; accept the same
object in the same snapshot. & Executed.\\
Reuse a proof after a rebuild & Refuse old epoch; reconstruct evidence in the new
context. & Executed.\\
Use positive length for \code{List Empty} & Admit constant-empty length law;
refute its \code{List Nat} version with a singleton. & Executed.\\
Forget equal-length correlation & Recompute the domain frame; do not reuse the
old request's result. & Executed.\\
Change coefficients, frame, or success label & Refuse certificate; accept the
unaltered certificate for the original request. & Executed.\\
Differentiate point equality & Reject that route; accept neighborhood equality
with the derivative evidence. & Written counterexample; native test pending.\\
Change measure beneath an a.e. fact & Require a new relation or checked transport;
accept the same measure's integral congruence. & Native test pending.\\
Drop a factor's module finiteness & Do not infer a finite representing algebra;
accept the finite-free family. & Written counterexample; native test pending.\\
Cast an inexact integer quotient & Require exactness and destination units;
accept the non-vacuous exact-quotient fixture. & Written arithmetic; native test pending.\\
\bottomrule
\end{tabularx}
\end{center}

Publication also needs tests for unused false intermediate assertions, changed
quantifiers, a wrong branch or precision, decorative use of a required method,
unlinked provider laws, and a correct proof of the wrong target. Passing the
small companion suite does not establish these document-level properties.

\subsection{Separate the library, services, type interface, and prose}
The study should compare the following arms, sharing mathematical packages
and implementation effort where appropriate:
\begin{center}\small
\begin{tabularx}{\textwidth}{@{}P{0.10\textwidth}Y@{}}
\toprule Arm & Capability allocation\\\midrule
$L_0$ & Ordinary Lean with its existing automation and idiomatic proof style.\\
$L_1$ & $L_0$ plus the same newly developed mathematical interfaces, model theorems,
and relevant theorem registrations supplied to Heather.\\
$L_2$ & $L_1$ plus the same added synthesis, certificate, and requirement-resolution
services, callable explicitly from ordinary Lean.\\
$L_3$ & $L_2$ plus Heather's automatic use-site proof role, fixed-meaning protocol,
shared evidence interface, and obligation explanations.\\
$L_4$ & $L_3$ plus the compact narrative input surface.\\
$R$ & A read-only mathematical renderer over the same checked Lean developments,
without a separate input language.\\
\bottomrule
\end{tabularx}
\end{center}

Ordinary Lean automatic parameters and reasonable thin wrappers are allowed
in the controls. A comparison that forbids those facilities would manufacture
an advantage for Heather. If a small Lean library reproduces the useful
interface, the correct conclusion is to ship that library rather than insist
that a new language has been validated.

The decisive comparisons are $L_3$ versus $L_2$ for the typing interface,
$L_4$ versus $L_3$ for narrative input, and $R$ versus the input-language arms
for reading benefits. Gains from a new mathematical theorem belong to the
library arm, not to syntax.

\subsection{Tasks, measurements, and cost accounting}
The inspected autonomous-derivative, tensor, and exact-quotient examples are
development tasks. They cannot become held-out evidence by being renamed.
After the interface and registrations are frozen, a maintainer should select
new related tasks and mutations without showing them to the implementation
team in advance. Both successful and unsuccessful tasks must be reported.

A counterbalanced study should use matched but distinct tasks so that a
participant's second attempt does not merely reuse the first proof from memory.
Record authoring time, interventions, failed repairs, proof-reading errors,
and success under a fixed assistance budget. Machine measurements should
separately record elaboration time, rule-generation time, rule firings, proof
term size, checking time, cache size, and rebuild cost. Kernel check time must
not be conflated with external search time.

Infrastructure costs include wrapper lemmas, registration metadata, provider
proofs, frame coverage theorems, reification bridges, training, and maintenance
after library changes. An amortized cost claim needs a stated number and kind
of future uses. A spectacularly short surface example with a bespoke package
behind it is not by itself a compression result.

The reading test should precede access to the expanded view. Readers should
recover the quantifier order, carrier and structures, domain restrictions,
result strength, and assumptions from the compact statement. An observed zero
error count is desirable but does not prove semantic perfection; adversarial
checking and formal interpretation work remain necessary.

\subsection{Answers to the second synthesis's ten implementation questions}
Table~\ref{tab:questions} uses the T1--T10 questions in \emph{Nine Refinements}
as a compact accountability ledger \cite{round3b}. The other synthesis's
concerns about source fidelity, non-vacuous fixtures, exact provider evidence,
realization, and equally equipped baselines are addressed by the same design
boundaries throughout the article \cite{round3a}.
\begin{center}\small
\begin{longtable}{@{}P{0.07\textwidth}P{0.32\textwidth}P{0.53\textwidth}@{}}
\caption{What Heather answers, and what remains to be measured.}\label{tab:questions}\\
\toprule Id & Question & Answer in this proposal\\\midrule\endfirsthead
\toprule Id & Question & Answer in this proposal\\\midrule\endhead
\bottomrule\endfoot
T1 & Cost of a real evidence index? & Not measured. The delivered trace is a toy;
native proof-size, timing, and snapshot data remain an acceptance gate.\\
T2 & Which rules fire, versus lookup? & One positivity implication fires in the
toy. A native registry must log use and compare against search and existing
property tactics.\\
T3 & Does simplification break anchors? & Only definitionally equal reuse,
checked equality/iff transport, or a registered relational consumer is allowed.
The toy refuses general rewriting.\\
T4 & Can reification be shared? & Share typed arithmetic syntax and interpretation
where structures and semantics agree; do not use one untyped expression model
for all domains. No general reifier was built here.\\
T5 & Completeness of Length abstractions? & Realizable frames give a written
criterion and four executed domain instances for the stated affine grammar,
not a theorem about every Leant Length contract.\\
T6 & How does a.e. reasoning enter? & As a relation indexed by its measure, with
an integral-congruence consumer and no automatic point-evaluation consumer.\\
T7 & Do lexical structures shorten preambles? & Unmeasured. Heather freezes
structures for fidelity; it does not claim to remove all instance setup.\\
T8 & What bounds rule generation? & A finite typed term pool, bounded constructors,
finite parameter instantiation, and explicit attempt limits precede closure.\\
T9 & What do compact-view readers recover? & A proposed blind reading protocol;
no study was run.\\
T10 & Is one implementation enough? & No independent implementation is delivered.
The Python producer/rechecker share code; a native checker and independent
adversarial implementation are still needed.\\
\end{longtable}
\end{center}

\section{Roadmap, stopping rules, and final recommendation}\label{sec:roadmap}

The next implementation should be a Lean package with a narrow public API,
not an ambitious parser. First, mechanize the list grammar interpretation and
the realizable-frame theorem, with the four admitted domains and exact source
correspondence. Export a theorem or a concrete refutation for the original
candidate. This is the behavioral slice for which the present article supplies
both a mathematical argument and executable development examples.

Second, add one property-implicit consumer over the already developed
exact-quotient interface. Resolve its designated proof parameters, retain its
fixed structures and source target, and measure it against equally equipped
Lean automatic parameters and explicit service calls. Include the satisfiable
fixture and the arithmetic mutations. A failed premise should remain a typed
obligation with its consuming operation.

Third, implement a small set of checked equality, proposition-equivalence,
and neighborhood adapters. Use the autonomous-derivative theorem as a
development test for preserving a quantified induction motive and range-local
hypotheses. Rebuild the index per declaration before attempting persistent
context migration. Then move to a genuinely new analysis task selected after
the implementation freezes.

Only after these gates should Heather add the narrative syntax illustrated in
Section~\ref{sec:surface}. A renderer can be developed earlier because it can
expose the same typed interface without creating a new input language.
A kernel primitive comes later still, and only if profiling identifies an
encoding cost the existing facilities cannot adequately remove.

The stopping rule is substantive. If $L_1$ explains the gains, retain the
mathematical library. If $L_2$ explains them, retain the services. If $R$ gives
the reading benefits, retain the renderer. If $L_3$ improves use-site authoring
but narrative input does not, ship the Lean extension without a separate
language. Each outcome preserves useful work; none requires claiming an
unobserved victory for Heather.

\subsection{Conclusion}
Heather treats concise mathematics as a problem of well-designed interfaces,
not of trusting omitted reasoning. Its objects carry proved local facets;
its operations demand precise evidence; its behavioral contracts refer to
actual functions and actual admissible inputs; its proof search cannot alter
the requested meaning. The ordinary Lean kernel remains the acceptance
boundary.

The main mathematical addition is a realizable-frame criterion that makes
finite affine reasoning sensitive to domains and correlations. The main
operational additions are concrete policies for proof-implicit arguments,
finite rule generation, rewriting, snapshot lifetime, and distinct failure
statuses. The delivered reference implementation exercises a small portion
of those policies and explicitly stops short of claiming a native elaborator
or Lean verification.

The question for the next stage is consequently testable: does making these
interfaces implicit and visible at the right moments reduce mathematical
proof-writing and repair costs beyond what the same theorems and services
already accomplish in Lean? Heather is worth pursuing to answer that question,
not because another name or a larger grammar settles it in advance.

\clearpage\appendix
\section{Source ledger and version boundaries}\label{app:sources}

The following full commit identifiers were returned by the repository interfaces
during this review. They identify the source observations used here, not an
assertion that the branches will remain at these revisions.
\par\noindent\begin{minipage}{\linewidth}
\begin{lstlisting}
Brainstorm
  bf3333905995060870f8585e297ffc16f3fe7e86
ProveIt
  24ce8bd743eaab64a91ce90725ea00f498d319d2
anthropics/fermats-last-theorem
  aa2d8b34692b16c70f699536de0d8e75b9a3e9ef
Leant
  3a40904be8a410d832d9ab6900b3d3e7b425eccb
\end{lstlisting}
\end{minipage}\par

The two requested main report sources are
\path{docs/round-3/synthesis/unified-report.tex} and
\path{docs/round-3/unified_report/unified_report.tex} in Brainstorm. Their
reconciled discussions, mathematical interfaces, source critiques, and
implementation questions are the basis of the inheritance analysis in this
article. This is not a claim of independent review of every proposal or every
separate appendix artifact referenced by those sources.

The focused Lean source study reads
\path{Analysis/FabiusFunction/Lean/FabiusFunction/AutonomousIteratedDeriv.lean}
in ProveIt and
\path{P2M/Sol/S_Algebra_exists_algHom_equiv_pi.lean}
in the Fermat repository. The focused Leant study reads
\path{src/Leant/Synth/Verification.hs},
\path{src/Leant/Synth/PostVerification.hs}, and
\path{src/Leant/Synth/Length/Adapter.hs}.

ProveIt's \path{lean-toolchain} at the observed pin specifies Lean 4.32.0.
The historical compilation claims in the syntheses concern their documented
runs and retained experiments. The current online Lean reference consulted
here identifies a different documentation version, 4.34.0-rc2. Heather's
uncompiled exports do not inherit a successful build result from either source.
The distinct Leant snapshot \code{823259f7} discussed by the reports is not
silently substituted for the \code{main} revision observed here.

\clearpage
\section{Minimal native interface contract}\label{app:interface}

A native implementation should start with a small semantic API. The following
is interface pseudocode, not compiled Lean syntax:
\par\noindent\begin{minipage}{\linewidth}
\begin{lstlisting}
EvidenceEntry:
  source_context : dependent telescope
  proposition    : elaborated Prop
  proof          : term checked at that proposition
  origin         : source assertion or registered application

Requirement:
  context        : current dependent telescope
  target         : fixed elaborated Prop
  consumer       : operation and designated parameter
  policy         : rules, providers, budgets, permitted axioms

Outcome:
  proved(proof_artifact)
  unresolved(frontier, exhaustion_information)
  refused(schema_or_context_mismatch)
  refuted(proof_of_negation)
\end{lstlisting}
\end{minipage}\par

A model counterexample belongs to a provider-specific intermediate result,
not to \code{refuted} until the source bridge and admissibility proof are
present. Request routing information and local snapshot identifiers accompany
these objects in the implementation, but do not substitute for any field above.

For the affine package, a mathematical native interface can use a base input
$s_0:S$, a finite family $s_j:S$, and a proof that every observation difference
is a rational linear combination of their differences. It also needs the
expression interpretation theorem and the correspondence to the actual
candidate. The simplest implementation should accept only explicitly
registered domain packages with those proofs. Automatic discovery of the
admissible image is not a prerequisite for the useful four instances in
Section~\ref{sec:frames}.

The native resolver should have a plain Lean invocation path in addition to
Heather's automatic proof-parameter path. This is both a practical escape hatch
and the strongest baseline for separating inference capability from use-site
convenience.

\clearpage
\section{Archive contents and reproduction}\label{app:reproduce}

\path{Heather.tex} is self-contained LaTeX source. It does not require downloaded
fonts, external figures, or a separate bibliography database.
\path{Heather.pdf} is its compiled rendering. The root README gives the build
commands.

The \path{companion/} directory contains the Python frontend, tests, use-site
demonstration, source examples, emitted result certificates, uncompiled Lean
candidates, and test logs. The status fields distinguish positive reference-model
conclusions from Lean acceptance. No \code{sorry}-based proof is presented as
verification; no successful Lean execution log is included.

To rerun the experiment from that directory:
\par\noindent\begin{minipage}{\linewidth}
\begin{lstlisting}
python test_heather.py
python use_site_demo.py
python heather.py examples/diagonal.hthr --output generated
python heather.py examples/empty.hthr --output generated
python heather.py examples/free.hthr --output generated
python heather.py examples/even.hthr --output generated
python heather.py examples/refuted.hthr --output generated
\end{lstlisting}
\end{minipage}\par
The test counts are deterministic; wall-clock timings in the JSON are not.
To attempt the exported Lean proofs separately, use a properly prepared Mathlib
project and invoke \code{lake env lean} on each generated file. Such an attempt
must be recorded as a new validation result. The package does not claim that
compilation in a new environment will be automatic or independent of its pin.

\clearpage
\addcontentsline{toc}{section}{References}
\begin{thebibliography}{99}\small

\bibitem{round3a}
\emph{From Properties to Proof Obligations: Round 3 Synthesis}.
Brainstorm campaign report, 6 September 2026.
\path{docs/round-3/synthesis/unified-report.tex}.
\href{https://github.com/VladimirReshetnikov/Brainstorm/blob/bf3333905995060870f8585e297ffc16f3fe7e86/docs/round-3/synthesis/unified-report.tex}{Pinned source}.
The source's historical experiments are attributed to the report, not rerun here.

\bibitem{round3b}
\emph{Nine Refinements: Property-Aware Typing in a Mathematician's Proof Language over Lean}.
Brainstorm campaign report, 6 September 2026.
\path{docs/round-3/unified_report/unified_report.tex}.
\href{https://github.com/VladimirReshetnikov/Brainstorm/blob/bf3333905995060870f8585e297ffc16f3fe7e86/docs/round-3/unified_report/unified_report.tex}{Pinned source}.

\bibitem{autonomous}
Vladimir Reshetnikov and ProveIt contributors.
\emph{Iterated derivatives of a solution of an autonomous equation}.
\path{AutonomousIteratedDeriv.lean}.
\href{https://github.com/VladimirReshetnikov/ProveIt/blob/24ce8bd743eaab64a91ce90725ea00f498d319d2/Analysis/FabiusFunction/Lean/FabiusFunction/AutonomousIteratedDeriv.lean}{Pinned source};
see \code{AutonomousODE.iteratedDeriv_eq_comp} and its differentiable-family corollary.

\bibitem{tensor}
Contributors to \emph{fermats-last-theorem}.
\path{P2M/Sol/S_Algebra_exists_algHom_equiv_pi.lean}.
\href{https://github.com/anthropics/fermats-last-theorem/blob/aa2d8b34692b16c70f699536de0d8e75b9a3e9ef/P2M/Sol/S_Algebra_exists_algHom_equiv_pi.lean}{Pinned source};
see the two-factor tensor lemma and \code{solution}.

\bibitem{leant-verification}
Vladimir Reshetnikov and Leant contributors.
\path{src/Leant/Synth/Verification.hs}.
\href{https://github.com/VladimirReshetnikov/Leant/blob/3a40904be8a410d832d9ab6900b3d3e7b425eccb/src/Leant/Synth/Verification.hs}{Pinned source};
see \code{Verified} and the candidate-group verification functions.

\bibitem{leant-post}
Vladimir Reshetnikov and Leant contributors.
\path{src/Leant/Synth/PostVerification.hs}.
\href{https://github.com/VladimirReshetnikov/Leant/blob/3a40904be8a410d832d9ab6900b3d3e7b425eccb/src/Leant/Synth/PostVerification.hs}{Pinned source};
see the epoch-scoped candidate handles and exact-permutation seal.

\bibitem{leant-adapter}
Vladimir Reshetnikov and Leant contributors.
\path{src/Leant/Synth/Length/Adapter.hs}.
\href{https://github.com/VladimirReshetnikov/Leant/blob/3a40904be8a410d832d9ab6900b3d3e7b425eccb/src/Leant/Synth/Length/Adapter.hs}{Pinned source};
see checked scalar and product-domain query preparation and its stated limits.

\bibitem{lean-reference}
Lean contributors. \emph{The Lean Language Reference}.
\href{https://lean-lang.org/doc/reference/latest/}{Online reference}, consulted
6 September 2026. The consulted landing page identifies version 4.34.0-rc2.

\bibitem{lean-application}
Lean contributors. \emph{The Lean Language Reference: Function Application}.
\href{https://lean-lang.org/doc/reference/latest/Terms/Function-Application/}{Official documentation},
consulted 6 September 2026. Includes automatic-parameter application behavior.

\bibitem{funprop}
Mathlib contributors. \emph{Mathlib.Tactic.FunProp}.
\href{https://leanprover-community.github.io/mathlib4_docs/Mathlib/Tactic/FunProp.html}{Generated primary API documentation},
consulted 6 September 2026.

\bibitem{mvcgen}
Lean contributors. \emph{The mvcgen Tactic: Overview}.
\href{https://lean-lang.org/doc/reference/latest/The--mvcgen--tactic/Overview/}{Official documentation},
consulted 6 September 2026.

\bibitem{lean-tactics}
Lean contributors. \emph{Tactic Reference} and \emph{The grind Tactic}.
\href{https://lean-lang.org/doc/reference/latest/Tactic-Proofs/Tactic-Reference/}{Tactic reference};
\href{https://lean-lang.org/doc/reference/latest/The--grind--tactic/}{grind documentation}.
Consulted 6 September 2026.

\bibitem{lean-axioms}
Lean contributors. \emph{The Lean Language Reference: Axioms}.
\href{https://lean-lang.org/doc/reference/latest/Axioms/}{Official documentation},
consulted 6 September 2026. Includes the distinction between proof axioms,
noncomputable data, parametricity, and trust-sensitive evaluation.

\bibitem{bochner}
Mathlib contributors. \emph{Mathlib.MeasureTheory.Integral.Bochner.Basic}.
\href{https://leanprover-community.github.io/mathlib4_docs/Mathlib/MeasureTheory/Integral/Bochner/Basic.html}{Generated primary API documentation},
consulted 6 September 2026; see \code{integral_congr_ae} and the integral definition.

\bibitem{liquid}
Patrick Rondon, Ming Kawaguchi, and Ranjit Jhala.
\emph{Liquid Types}. PLDI, 2008.
\href{https://goto.ucsd.edu/~rjhala/papers/liquid_types.html}{Authors' publication page and abstract}.

\bibitem{occurrence}
Sam Tobin-Hochstadt and Matthias Felleisen.
\emph{The Design and Implementation of Typed Scheme: From Scripts to Programs}.
ArXiv preprint 1106.2575, 2011.
\href{https://arxiv.org/abs/1106.2575}{Authors' preprint abstract}.

\end{thebibliography}
\end{document}
